\chapter{Experiments and Results}

\section{Evaluation Setting}

Initially, I intended to evaluate the proposed method on the UAV dataset introduced by AirPose \cite{saini2022airpose}, as the authors provide calibration matrices for the two drones used during acquisition. In principle, this makes 
it possible to reconstruct 3D joint coordinates through triangulation and use them as pseudo ground truth for quantitative evaluation.

However, in practice, the triangulated 3D poses turned out to be excessively noisy. The reconstructed joints were highly unstable due to synchronization inaccuracies, imperfect detections, and sensitivity of triangulation to small 
2D localization errors. As a consequence, the resulting 3D annotations were not reliable enough to produce meaningful quantitative comparisons. The evaluation metrics computed on these reconstructions were inconsistent and did not 
reflect the true quality of the predicted poses.

For this reason, I performed the quantitative evaluation on the test split of AthletePose3D \cite{yeung2025athletepose3d}, where accurate 3D ground-truth annotations are available. This allows a reliable and reproducible comparison 
between models under controlled conditions.

\section{Evaluation Protocol}

During evaluation, 2D keypoints are first extracted and normalized following the preprocessing procedure described in the Section~\ref{sec:preprocess}. Temporal windows of 81 consecutive frames are fed to the network using a 
sliding-window strategy consistent with training.

Predicted 3D poses are root-centered before metric computation. When enabled, test-time augmentation (TTA) through horizontal flipping is applied, and final predictions are obtained by averaging the outputs of the original and 
flipped inputs.

\section{Evaluation Metrics}

To assess performance, I report Mean Per Joint Position Error (MPJPE) and Procrustes-aligned MPJPE (P-MPJPE), both measured in millimeters.

MPJPE computes the average Euclidean distance between predicted and ground-truth 3D joint coordinates. P-MPJPE applies a rigid Procrustes alignment before error computation, removing global rotation and translation differences. 
This metric isolates structural pose accuracy independently of global orientation.


\section{Quantitative Results}

\begin{table}[H]
\centering
\caption{Quantitative comparison performed on the AthletePose3D test set.}
\label{tab:results}
\begin{tabular}{lcc}
\hline
Model & MPJPE (mm) & P-MPJPE (mm) \\
\hline
MotionBERT & -- & -- \\
MotionAGFormer & 160.6560 & 86.1506 \\
MotionAGFormer Finetuned & 14.1359 & 11.0060 \\
Baseline & 45.2764 & 4.3332 \\
Brancher & 30.9958 & 27.8366 \\
Brancher Self Supervised Finetuned & -- & -- \\
\hline
\end{tabular}
\end{table}

Table~\ref{tab:results} reports the quantitative performance of all evaluated models on the AthletePose3D test split.

The original MotionAGFormer model \cite{mehraban2024motionagformer}, trained on Human3.6M datasets \cite{zhu2023motionbert}, exhibits the highest error, with an MPJPE of 160.66 mm and a P-MPJPE of 86.15 mm. After fine-tuning on 
AthletePose3D, MotionAGFormer Finetuned achieves a substantial reduction in error, reaching 14.14 mm MPJPE and 11.01 mm P-MPJPE.

The Baseline model obtains an MPJPE of 45.28 mm and a P-MPJPE of 4.33 mm. Compared to MotionAGFormer Finetuned, it presents a higher MPJPE but a lower Procrustes-aligned error.

The proposed Brancher architecture achieves an MPJPE of 31.00 mm and a P-MPJPE of 27.84 mm, improving over the Baseline in terms of MPJPE while yielding a higher P-MPJPE.

TODO: da finire e dire che questa era una prima analisi dei risultati basata unicamennte sui numeri.


\section{Qualitative Analysis}
